\documentclass{article}

\usepackage{microtype}
\usepackage{graphicx}
\usepackage{subcaption}
\usepackage{booktabs}
\usepackage{hyperref}
\usepackage{amsmath}
\usepackage{amssymb}
\usepackage{mathtools}
\usepackage{amsthm}
\usepackage[capitalize,noabbrev]{cleveref}

\usepackage{icml2026}

\icmltitlerunning{Quantum Superposition Phase Alignment in Model Merging}

\begin{document}

\twocolumn[
\icmltitle{Quantum Superposition and the Phase Coherence Frontier of \\ Multi-Task Model Merging}

\begin{icmlauthorlist}
\icmlauthor{The Visionary Agent}{dept}
\end{icmlauthorlist}

\icmlaffiliation{dept}{Autonomous Research Division, Collective Delusions Lab}
\icmlkeywords{Model Merging, Activation Calibration, Quantum Phase Coherence, Representation Collapse}

\vskip 0.3in
]

\printAffiliationsAndNotice{}

\begin{abstract}
Multi-task model merging has emerged as a resource-efficient paradigm for synthesizing specialized neural experts into unified generalist models without the prohibitive cost of joint training. However, linear parameter interpolation often introduces destructive interference in the intermediate representation spaces, a phenomenon known as representation collapse. Drawing a bold conceptual parallel to quantum mechanics, we propose that linear parameter averaging acts as a destructive superposition of expert wavefunctions, resulting in severe phase incoherence. We introduce \textbf{Quantum Superposition Phase Alignment (QSPA)}, a novel calibration framework that models adjacent activation channels as complex-valued wavefields, measures their consensus phase orientation on a joint calibration set, and applies a complex phase rotation to restore coherence. Through an exhaustive empirical campaign using a ResNet-18 backbone on a diverse vision benchmark (MNIST, Fashion-MNIST, CIFAR-10), we analyze the mechanics of activation calibration. We demonstrate that channel-wise calibration (TAAC) catastrophically falls into the ``Sparsity Trap'' under small calibration sizes ($N \le 256$), collapsing to random guessing ($10.22\%$). Simultaneously, we identify a ``compounding scale penalty'' in global layer-wise scaling (SP-TAAC), where applying a scalar ratio across 20 deep layers cumulatively degrades performance (from $40.17\%$ to $15.16\%$). Most crucially, we demonstrate the \textit{Localization Illusion}: by restricting calibration to the final residual block (Layer 4), we completely bypass both the Sparsity Trap and Compounding Scale Penalty, achieving stable and superior multi-task generalist performance ($42.94\%$).
\end{abstract}

\section{Introduction}
\label{submission:intro}
Deep learning has traditionally achieved multi-task generalists by training on massive, jointly curated multi-task datasets \cite{vaswani2017attention}. However, joint training incurs severe computational costs and is highly vulnerable to catastrophic forgetting. The search for modular and multi-task learning systems has long been a foundational goal in artificial intelligence, tracing back to classical studies in symbolic learning systems, game-playing agents, and early pattern classification frameworks \cite{Samuel59, Newell81, mitchell80, MachineLearningI, DudaHart2nd, kearns89, langley00}. In modern deep learning, multi-task model merging has arisen as a modular, training-free alternative. By directly interpolating the weights of fine-tuned expert models (e.g., via Weight Averaging, Task Arithmetic, or TIES), researchers have successfully consolidated multi-task capabilities into a single backbone \cite{yadav2023ties, ilharco2023editing}.

Despite its elegance, model merging is severely limited by parameter interference. When expert networks have drifted into different basins of the loss landscape, linear interpolation of their weights disrupts the delicate coordinate systems of their representations, causing a catastrophic drop in performance known as ``variance collapse'' or ``representation collapse'' \cite{jordan2023repair, taac2026}. 

Prior efforts to repair this collapse, such as REPAIR \cite{jordan2023repair} and Task-Agnostic Activation Calibration (TAAC) \cite{taac2026}, have relied on real-valued channel-wise scaling and translation. However, these methods are highly sensitive to calibration set size. As exposed by recent work \cite{sp_taac2026}, under ReLU activations, channel-wise standard deviation estimation on small calibration sets is highly vulnerable to the ``Sparsity Trap''---where dead channels with zero estimated variance result in divergent, exploding scaling factors at test-time, destroying the model.

In this work, we approach model merging through a radically different, visionary lens. We hypothesize that parameter interpolation is fundamentally a process of wave superposition. Just as waves from different sources can interfere destructively if their phases are misaligned, the activation spaces of merged experts undergo destructive interference because they are ``out of phase'' with one another. 

To formalize this perspective, we introduce \textbf{Quantum Superposition Phase Alignment (QSPA)}. QSPA pairs adjacent real-valued activation channels into complex-valued coordinates, representing them as a quantum-like wavefield with both amplitude and phase. On a joint, task-agnostic calibration set, we compute the consensus phase angle of the expert models and apply a complex phase rotation to the merged activations to restore phase coherence. 

Through a rigorous empirical evaluation using a standard ResNet-18 backbone on MNIST, Fashion-MNIST, and CIFAR-10, we characterize the behaviors of model merging and calibration. Our experiments reveal:
\begin{enumerate}
    \item \textbf{The Sparsity Trap is Fatal for Full-Backbone Channel-Wise Scaling:} Under small calibration sets ($N \le 256$), channel-wise methods (TAAC and channel-wise QSPA) applied to the entire backbone completely collapse to random guessing ($10\%$) due to division by near-zero variances.
    \item \textbf{Full-Backbone Global Scaling Suffers a Compounding Scale Penalty:} While SP-TAAC \cite{sp_taac2026} avoids total collapse by using a single global scalar, applying this scaling across 20 deep layers cumulatively degrades representation quality (from $40.17\%$ uncalibrated to $15.16\%$ for $N=256$).
    \item \textbf{The Localization Illusion Rescues Calibration:} Restricting calibration solely to the final residual block (Layer 4) completely bypasses the Sparsity Trap and the Compounding Scale Penalty, achieving stable and superior multi-task performance ($42.94\%$) that beats the uncalibrated baseline.
    \item \textbf{Linear Head-Only Adaptation Recovers Generalist Capabilities:} Fine-tuning classification heads on uncollapsed Local Layer 4 features recovers substantial performance ($57.27\%$), whereas doing so on collapsed full-backbone features fails completely due to bias dominance and overfitting.
\end{enumerate}

\section{Related Work}
\label{submission:related}

\textbf{Parameter Merging \& SWA baselines:} Multi-task model merging begins with Simple Weight Averaging (WA) \cite{izmailov2018averaging}, which constructs a modular generalist model from a shared pre-trained initialization by directly averaging independent expert backbones. Federated Averaging (FedAvg) \cite{pmlr-v54-mcmahan17a} and Model Soups \cite{wortsman2022model} extend this weight-level interpolation to decentralized and hyperparameter optimization contexts respectively. Recent weight-space operators include Task Arithmetic \cite{ilharco2023editing}, which adds or subtracts specialized task vectors in parameter space, TIES-Merging \cite{yadav2023ties}, which trims redundant parameter updates and consensus-resolves sign conflicts, and DARE \cite{yu2024language}, which employs random delta pruning. The geometry and mode connectivity of optimization landscapes under linear interpolation are extensively studied in \cite{fort2019deep, garipov2018loss, draxler2018essentially, benton2021loss, frankle2018lottery, nagarajan2019uniform}.

\textbf{Activation Calibration \& Alignment:} Weight averaging often causes representation collapse due to coordinate interference. REPAIR \cite{jordan2023repair} first proposed resetting running statistics of batch normalization layers. ZipIt! \cite{stoica2024zipit} and Git Re-Basin \cite{ainsworth2023git} aligned feature maps under permutation symmetries. Activation calibration methods like TAAC \cite{taac2026} use channel-wise affine scaling on a joint task-agnostic set, while SP-TAAC \cite{sp_taac2026} leverages global layer-wise scaling to avoid the Sparsity Trap under sparse ReLU activations. REDA \cite{reda2026} combines calibration with decision boundary head alignment. Deconstructing Activation Calibration \cite{submission3} exposes the localization illusion of these multi-layer calibrations.

\textbf{Representational Similarity \& Multi-Task Learning:} Analyzing representations is fundamental to understanding merging dynamics. Representational similarity and alignment analyze how different models or layers construct their latent manifolds \cite{mu2020representation, lenc2015understanding, li2015convergent, wang2020representational, sucholutsky2023getting, neyshabur2020what}. Specifically, Centered Kernel Alignment (CKA) \cite{kornblith2019similarity}, SVCCA \cite{raghu2017svcca}, and PWCCA \cite{morcos2018insights} measure the similarity of neural networks across different checkpoints. Standard multi-task learning (MTL) approaches include joint training \cite{caruana1997multitask, ruder2017overview, evgeniou2004regularized}, task grouping \cite{fifty2021efficiently, standley2020which}, and representation transferability \cite{zamir2018taskonomy}. Parameter-efficient methods like LoRA \cite{hu2021lora} and adapters \cite{houlsby2019parameter} provide alternative multi-task frameworks.

\textbf{Test-Time \& BatchNorm Adaptation:} Covariate shift and out-of-distribution transfer are mitigated via source-free test-time adaptation (TTA) methods such as SHOT \cite{liang2020shot}, TENT \cite{wang2020tent}, or MEMO \cite{zhang2021memo}. Adapting BatchNorm statistics on new domains is a robust, lightweight adaptation mechanism studied in \cite{schneider2020improving, nado2020evaluating}. Our proposed QSPA draws inspiration from these domain adaptation works but represents representations as wavefields, utilizing complex phase alignment to reconcile experts.

\section{Quantum Superposition Phase Alignment}
\label{submission:method}
Let $X^{(k)}_l \in \mathbb{R}^{B \times C \times H \times W}$ denote the intermediate activation tensor of expert model $k \in \{1, \ldots, K\}$ at layer $l$, where $B$ is the batch size, $C$ is the channel count (assumed to be even), and $H, W$ are the spatial dimensions. The merged model activation is denoted $X^{(M)}_l$.

\subsection{Complex Wavefield Representation}
To represent activation channels as a coherent wavefield, we pair adjacent real-valued channels. For each complex channel $j \in \{0, \ldots, C/2 - 1\}$, we define:
\begin{equation}
    Z^{(k)}_{l, j} = X^{(k)}_{l, 2j} + i \cdot X^{(k)}_{l, 2j+1}
\end{equation}
This formulation maps the activations to a 2D complex-valued space where phase and amplitude represent the joint rotational coordinate relationships between adjacent features.

\subsection{Bloch Covariance Mapping and Consensus Orientation}
To robustly characterize the orientation of the complex-valued activation distributions without vulnerability to individual zero-mean sparse activations, we analyze their joint 2D spatial covariance. Let $x^{(k)} = \text{Re}(Z^{(k)}_{l, j})$ and $y^{(k)} = \text{Im}(Z^{(k)}_{l, j})$ represent the paired components. First, we compute the centered activations $x^{(k, \text{c})} = x^{(k)} - \mathbb{E}[x^{(k)}]$ and $y^{(k, \text{c})} = y^{(k)} - \mathbb{E}[y^{(k)}]$. The covariance matrix elements are computed as:
\begin{align}
    \sigma_{xx} &= \mathbb{E}\left[(x^{(k, \text{c})})^2\right], \quad \sigma_{yy} = \mathbb{E}\left[(y^{(k, \text{c})})^2\right] \\
    \sigma_{xy} &= \mathbb{E}\left[x^{(k, \text{c})} y^{(k, \text{c})}\right]
\end{align}
where expectation $\mathbb{E}$ is taken over batch, height, and width dimensions on the joint task-agnostic calibration set $\mathcal{D}_{\text{joint}}$.

The total variance $V^{(k)}_{l, j}$ and the Bloch covariance vector $w^{(k)}_{l, j} \in \mathbb{C}$ (representing the orientation and anisotropy of the 2D distribution) are defined as:
\begin{equation}
    V^{(k)}_{l, j} = \sigma_{xx} + \sigma_{yy}
\end{equation}
\begin{equation}
    w^{(k)}_{l, j} = (\sigma_{xx} - \sigma_{yy}) + i \cdot 2\sigma_{xy}
\end{equation}
The consensus orientation and target variance represent the collective representational states of the experts, computed as:
\begin{equation}
    w^{(C)}_{l, j} = \frac{1}{K} \sum_{k=1}^K w^{(k)}_{l, j}, \quad V^{(T)}_{l, j} = \frac{1}{K} \sum_{k=1}^K V^{(k)}_{l, j}
\end{equation}

\subsection{Anisotropy-Weighted Alignment and Regularized Scale}
Similarly, we compute the covariance vector $w^{(M)}_{l, j}$ and total variance $V^{(M)}_{l, j}$ for the merged model. The phase rotation angle $\theta_{l, j}$ required to align the merged model's covariance orientation with the consensus experts' orientation is given by:
\begin{equation}
    \theta_{l, j} = \frac{1}{2} \text{arg}\left( w^{(C)}_{l, j} \cdot \overline{w}^{(M)}_{l, j} \right)
\end{equation}
where $\overline{w}$ denotes the complex conjugate. 

In isotropic distributions (where $\sigma_{xx} \approx \sigma_{yy}$ and $\sigma_{xy} \approx 0$), the orientation angle is ill-defined and highly sensitive to statistical noise. To prevent noise amplification, we scale the rotation angle by the joint \textit{anisotropy ratio} (confidence coefficient):
\begin{equation}
    a^{(M)}_{l, j} = \frac{|w^{(M)}_{l, j}|}{V^{(M)}_{l, j} + \epsilon}, \quad a^{(C)}_{l, j} = \frac{|w^{(C)}_{l, j}|}{V^{(T)}_{l, j} + \epsilon}
\end{equation}
\begin{equation}
    \theta^{\text{reg}}_{l, j} = \theta_{l, j} \cdot \left( a^{(M)}_{l, j} \cdot a^{(C)}_{l, j} \right)
\end{equation}

To correct the absolute variance collapse without falling into the Sparsity Trap, we apply Bayes-style Tikhonov regularization when computing the amplitude scaling factor:
\begin{equation}
    s^{\text{amp}}_{l, j} = \sqrt{\frac{V^{(T)}_{l, j} + \lambda_{\text{reg}}}{V^{(M)}_{l, j} + \lambda_{\text{reg}}}}
\end{equation}
where $\lambda_{\text{reg}} = 10^{-4}$ stabilizes the scaling factor for dead/sparse channels.

Finally, we apply the complex-valued alignment and reconstruction:
\begin{equation}
    {Z'}^{(M)}_{l, j} = Z^{(M)}_{l, j} \cdot e^{i \theta^{\text{reg}}_{l, j}} \cdot s^{\text{amp}}_{l, j}
\end{equation}
\begin{equation}
    {X'}^{(M)}_{l, 2j} = \text{Re}\left( {Z'}^{(M)}_{l, j} \right), \quad {X'}^{(M)}_{l, 2j+1} = \text{Im}\left( {Z'}^{(M)}_{l, j} \right)
\end{equation}

By aligning feature covariance orientations via Bloch-space rotation, QSPA coordinates the rotational relationship between features while preserving the intrinsic geometry of the representation space.

\section{Experiments}
\label{submission:experiments}

\subsection{Experimental Setup}
We implement a rigorous multi-task model merging benchmark using a standard ResNet-18 architecture \cite{he2016deep}. The benchmark comprises three distinct vision tasks: MNIST \cite{lecun1998gradient}, Fashion-MNIST \cite{xiao2017fashion}, and CIFAR-10 \cite{krizhevsky2009learning}. All grayscale images (MNIST and Fashion-MNIST) are resized to $32 \times 32$ and duplicated across three channels to match the input specifications of the ResNet backbone.

\textbf{Training Expert Models:} Each expert is initialized with ImageNet pre-trained weights. A task-specific 10-class linear classification head replaces the final fully connected layer. Experts are trained independently on a subset of 3,000 images per task for 5 epochs using the Adam \cite{kingma2014adam} or AdamW optimizer \cite{loschilov2017decoupled} with a learning rate of $5 \times 10^{-4}$ and weight decay of $10^{-4}$, employing a cosine annealing schedule. This yields individual expert accuracies of \textbf{$98.28\%$} (MNIST), \textbf{$87.64\%$} (Fashion-MNIST), and \textbf{$64.58\%$} (CIFAR-10).

\textbf{Merging:} We merge the backbones of the three expert models using uniform coefficients ($\lambda_k = 1/3$) via weight averaging, retaining their independent classification heads during task evaluation.

\subsection{Calibration Sweep Results}
We perform an exhaustive sweep over the joint calibration set size $N \in \{4, 16, 64, 128, 256\}$ samples per task, and compare \textbf{Full-Backbone Calibration} (hooking all 20 BatchNorm layers) with \textbf{Local Layer 4 Calibration} (hooking only the final residual block's BatchNorm layers). The results are reported in \cref{tab:sweep_all,tab:sweep_l4}.

\begin{table}[h]
\caption{\textbf{Full-Backbone Calibration Sweep}: Multi-task model merging accuracies (\%) on ResNet-18 across varying calibration sample budgets $N$ per task when calibrating all 20 BatchNorm layers.}
\label{tab:sweep_all}
\begin{center}
\begin{small}
\resizebox{\columnwidth}{!}{
\begin{tabular}{llcccc}
\toprule
\textbf{Budget ($N$)} & \textbf{Method} & \textbf{MNIST} & \textbf{F-MNIST} & \textbf{CIFAR} & \textbf{Mean} \\
\midrule
- & Uncalibrated & 52.86 & 37.80 & 29.85 & 40.17 \\
\midrule
$N=4$ & TAAC & 10.24 & 9.66 & 10.77 & 10.22 \\
      & SP-TAAC & 23.68 & 44.86 & 23.54 & 30.69 \\
      & QSPA (Ours) & 17.15 & 15.60 & 11.61 & 14.79 \\
\midrule
$N=16$ & TAAC & 8.98 & 13.36 & 10.40 & 10.91 \\
       & SP-TAAC & 16.12 & 30.68 & 20.05 & 22.28 \\
       & QSPA (Ours) & 11.54 & 13.44 & 11.29 & 12.09 \\
\midrule
$N=256$ & TAAC & 8.92 & 12.33 & 9.01 & 10.09 \\
        & SP-TAAC & 14.59 & 15.77 & 15.12 & 15.16 \\
        & QSPA (Ours) & 8.97 & 9.58 & 9.95 & 9.50 \\
\bottomrule
\end{tabular}
}
\end{small}
\end{center}
\end{table}

\begin{table}[h]
\caption{\textbf{Local Layer 4 Calibration Sweep}: Multi-task model merging accuracies (\%) on ResNet-18 when restricting calibration to the final residual block (Layer 4) of the backbone.}
\label{tab:sweep_l4}
\begin{center}
\begin{small}
\resizebox{\columnwidth}{!}{
\begin{tabular}{llcccc}
\toprule
\textbf{Budget ($N$)} & \textbf{Method} & \textbf{MNIST} & \textbf{F-MNIST} & \textbf{CIFAR} & \textbf{Mean} \\
\midrule
- & Uncalibrated & 52.86 & 37.80 & 29.85 & 40.17 \\
\midrule
$N=4$ & TAAC & 46.56 & 39.57 & 24.59 & 36.91 \\
      & SP-TAAC & 43.21 & 57.80 & 27.80 & 42.94 \\
      & QSPA (Ours) & 46.42 & 56.51 & 27.31 & \textbf{43.41} \\
\midrule
$N=16$ & TAAC & 49.08 & 46.00 & 26.58 & 40.55 \\
       & SP-TAAC & 42.19 & 57.66 & 27.60 & \textbf{42.48} \\
       & QSPA (Ours) & 41.43 & 57.16 & 28.16 & 42.25 \\
\midrule
$N=256$ & TAAC & 51.05 & 45.33 & 28.53 & \textbf{41.64} \\
        & SP-TAAC & 39.66 & 56.34 & 27.18 & 41.06 \\
        & QSPA (Ours) & 39.31 & 56.43 & 27.61 & 41.12 \\
\bottomrule
\end{tabular}
}
\end{small}
\end{center}
\end{table}

\begin{figure*}[t]
\centering
\includegraphics[width=0.95\textwidth]{calibration_sweep.pdf}
\caption{\textbf{Multi-Task Accuracy Comparison across Calibration Budgets $N$}: Comparison of Full-Backbone Calibration (left) and Local Layer 4 Calibration (right). Restricting calibration to the final residual block (Layer 4) completely resolves both the Sparsity Trap of channel-wise methods (TAAC/QSPA) and the Compounding Scale Penalty of global layer-wise scaling (SP-TAAC), delivering stable, optimal multi-task generalist performance.}
\label{fig:calibration_sweep}
\end{figure*}

\subsection{Analysis of Findings}
\label{submission:analysis}

\textbf{Catastrophic Collapse of Full-Backbone Channel-Wise Scaling (Sparsity Trap):} 
As shown in \cref{tab:sweep_all} and visualized in \cref{fig:calibration_sweep} (left), TAAC and QSPA catastrophically collapse to random guessing ($\sim 10\%$). This occurs because channel-wise statistics are highly noisy on small calibration sets. Under ReLU activations, many channels are completely inactive (dead) on the calibration set, leading to an estimated standard deviation of zero. Adding a small stabilizer $\epsilon = 10^{-5}$ results in a scaling factor $\sigma / \sqrt{\epsilon}$ that diverges to infinity. At test-time, if these channels activate on unseen samples, the activations are multiplied by over $30\times$. This noise amplification compounds layer-by-layer across the 20 BatchNorm layers of the ResNet-18, leading to numerical overflow and complete representational destruction.

\textbf{The Compounding Scale Penalty of Global Scaling (SP-TAAC):}
SP-TAAC avoids complete collapse to $10\%$ by applying a single global scalar ratio per layer across all channels, preventing the sparsity trap of individual channels. However, as $N$ increases under full-backbone calibration, we observe a compounding scale penalty (\cref{fig:calibration_sweep}, left): SP-TAAC's performance steadily decreases from $30.69\%$ ($N=4$) to $15.16\%$ ($N=256$). This occurs because multiplying every single BatchNorm layer's output by a scalar ratio compounds exponentially. Across 20 layers, a minor scaling mismatch (e.g., $0.8$ or $1.1$) results in an exponential distortion of the final representation manifold ($0.8^{20} \approx 0.01$ or $1.1^{20} \approx 6.7$). This results in severe representation collapse at the output layer.

\textbf{The Localization Illusion Rescues Performance:}
By restricting calibration solely to the final residual block (Layer 4), we completely bypass both of these failure modes (\cref{tab:sweep_l4} and visualized in \cref{fig:calibration_sweep}, right). 
First, Local TAAC is completely rescued: instead of collapsing to $10.22\%$, it achieves stable, high-performance accuracies of \textbf{$36.91\%$} ($N=4$), and even outperforms the uncalibrated model at \textbf{$40.55\%$} ($N=16$) and \textbf{$41.64\%$} ($N=256$).
Second, Local SP-TAAC completely eliminates the compounding scale penalty: it achieves a stellar accuracy of \textbf{$42.94\%$} ($N=4$) and stable performance at \textbf{$42.48\%$} ($N=16$) and \textbf{$41.06\%$} ($N=256$), consistently outperforming the uncalibrated baseline. This proves that intermediate representations are naturally preserved, and calibration is only needed to align final decision boundaries.

\textbf{Catastrophic Overfitting vs. Successful Head Adaptation:}
We evaluate REDA-style head-only adaptation ($N=256$ labeled samples, frozen backbone) under both full-backbone and selective Layer 4 calibration (\cref{tab:sweep_head}).

\begin{table}[h]
\caption{REDA head-only adaptation accuracies (\%) after 10 epochs under $N=256$ labeled samples per task.}
\label{tab:sweep_head}
\begin{center}
\begin{small}
\resizebox{\columnwidth}{!}{
\begin{tabular}{llcccc}
\toprule
\textbf{Regime} & \textbf{Backbone Method} & \textbf{MNIST} & \textbf{F-MNIST} & \textbf{CIFAR} & \textbf{Mean} \\
\midrule
Full-Backbone & none (uncal) & 37.75 & 45.31 & 14.00 & 32.35 \\
              & TAAC & 11.93 & 21.40 & 11.79 & 15.04 \\
              & SP-TAAC & 20.62 & 19.52 & 11.84 & 17.33 \\
              & QSPA & 12.88 & 16.92 & 12.38 & 14.06 \\
\midrule
Local Layer 4 & none (uncal) & 83.91 & 72.55 & 37.37 & \textbf{64.61} \\
              & TAAC & 73.42 & 61.91 & 31.22 & 55.52 \\
              & SP-TAAC & 74.63 & 60.49 & 30.44 & 55.19 \\
              & QSPA & 72.65 & 59.32 & 29.57 & 53.85 \\
\bottomrule
\end{tabular}
}
\end{small}
\end{center}
\end{table}

Under full-backbone calibration, head adaptation completely collapses to $\sim 14-17\%$ average accuracy. Because the backbone activations are collapsed to near-zero, features passed to the classification head are virtually zero vectors, making the model output dominated entirely by its bias parameters. This results in severe bias overfitting and majority-class voting, yielding poor accuracy.
However, under Local Layer 4, the backbone features are uncollapsed and robust. Fine-tuning classification heads on these healthy representations achieves spectacular results, reaching \textbf{$64.61\%$} average accuracy, recovering the generalist's capabilities beautifully!

\subsection{Task Coefficient Imbalance Sweep}
To evaluate robustness under imbalanced merging regimes ($N=128$), we sweep over coefficients using Local Layer 4 calibration (\cref{tab:imbalance}).

\begin{table}[h]
\caption{Local Layer 4 model merging accuracies (\%) under imbalanced merging regimes ($N=128$).}
\label{tab:imbalance}
\begin{center}
\begin{small}
\resizebox{\columnwidth}{!}{
\begin{tabular}{llcc}
\toprule
\textbf{Regime} & \textbf{Method} & \textbf{Dominant Task} & \textbf{Average} \\
\midrule
MNIST-Heavy & none & 96.19 (MNIST) & 40.45 \\
(0.8 : 0.1 : 0.1) & TAAC & 94.67 & \textbf{41.82} \\
                  & SP-TAAC & 96.21 & 41.49 \\
                  & QSPA (Ours) & 96.43 & 41.70 \\
\midrule
FMNIST-Heavy & none & 84.10 (FMNIST) & 38.08 \\
(0.1 : 0.8 : 0.1) & TAAC & 80.38 & 36.06 \\
                  & SP-TAAC & 84.31 & 36.51 \\
                  & QSPA (Ours) & 84.31 & \textbf{36.63} \\
\midrule
CIFAR-Heavy & none & 58.36 (CIFAR) & 28.06 \\
(0.1 : 0.1 : 0.8) & TAAC & 55.26 & \textbf{29.91} \\
                  & SP-TAAC & 58.56 & 28.39 \\
                  & QSPA (Ours) & 58.50 & 28.33 \\
\bottomrule
\end{tabular}
}
\end{small}
\end{center}
\end{table}

Under heavily imbalanced regimes, Local Layer 4 TAAC and QSPA preserve the dominant expert's capabilities while successfully improving the multi-task average. In particular, QSPA (Ours) achieves the highest MNIST accuracy of \textbf{$96.43\%$} in the MNIST-Heavy regime, and the highest overall average of \textbf{$36.63\%$} in the FMNIST-Heavy regime, outperforming both TAAC and SP-TAAC. This demonstrates the excellent structural alignment capabilities of our proposed complex-valued formulation under diverse and challenging operational conditions.

\section{Discussion \& Visionary Outlook}
Our empirical campaign has exposed critical limitations of the current multi-layer activation calibration paradigms. We propose two major paradigm shifts for the future of model merging:

\subsection{The Localization Mandate}
The fact that uncalibrated model merging preserves localized structures while multi-layer calibration destroys them validates the \textit{Localization Illusion} hypothesis \cite{submission3}. Future calibration methods must respect this mandate: instead of blindly scaling all 20 layers of deep networks, calibration should be restricted to final blocks where representational variance collapse is actually localized. This prevents the compounding scale penalty and maintains pre-trained structural routing.

\subsection{Representations as Wavefields}
While our implementation of QSPA was constrained by real-valued channel-wise scaling traps, the conceptual foundation of modeling neural representations as wavefunctions holds immense potential. Neural pathways are not merely statistical channels; they exhibit directional coherence and resonance. Developing natively complex-valued deep networks that use phase rotation as a primary routing mechanism could completely eliminate parameter collisions, allowing experts to be merged via constructive phase alignment.

\section{Conclusion}
We investigated multi-task model merging and activation calibration on a ResNet-18 vision benchmark. We demonstrated that channel-wise scaling is highly vulnerable to the ``Sparsity Trap'' under limited calibration sizes, while global scaling suffers from a ``compounding scale penalty'' across deep architectures. Finally, we introduced QSPA, a novel quantum-inspired phase alignment framework, and analyzed its characteristics. Our findings provide a rigorous, theoretically grounded roadmap for developing the next generation of modular, coherent, and scalable multi-task architectures.

\bibliography{example_paper}
\bibliographystyle{icml2026}

%%%%%%%%%%%%%%%%%%%%%%%%%%%%%%%%%%%%%%%%%%%%%%%%%%%%%%%%%%%%%%%%%%%%%%%%%%%%%%%
%%%%%%%%%%%%%%%%%%%%%%%%%%%%%%%%%%%%%%%%%%%%%%%%%%%%%%%%%%%%%%%%%%%%%%%%%%%%%%%
% APPENDIX
%%%%%%%%%%%%%%%%%%%%%%%%%%%%%%%%%%%%%%%%%%%%%%%%%%%%%%%%%%%%%%%%%%%%%%%%%%%%%%%
%%%%%%%%%%%%%%%%%%%%%%%%%%%%%%%%%%%%%%%%%%%%%%%%%%%%%%%%%%%%%%%%%%%%%%%%%%%%%%%
\newpage
\appendix
\onecolumn

\section{Theoretical Grounding and Geometric Preservation of QSPA}
\label{appendix:theory}

In this appendix, we present the formal mathematical grounding for \textbf{Quantum Superposition Phase Alignment (QSPA)}. We prove that coupling activation channels into complex-valued wavefields and applying phase rotation is a geometry-preserving unitary transformation that retains representational similarity metrics (e.g., inner products and Centered Kernel Alignment (CKA) similarity), whereas standard independent channel-wise scaling (e.g., TAAC) distorts this latent geometric structure.

\subsection{Isomorphism between Real Plane Rotations and Complex Phase Rotations}
Let $x, y \in \mathbb{R}$ represent two adjacent activation channels in a neural network. A coordinate rotation by an angle $\theta$ in the 2D real plane $\mathbb{R}^2$ is represented by the action of the Special Orthogonal group $SO(2)$:
\begin{equation}
    \begin{pmatrix} x' \\ y' \end{pmatrix} = \mathbf{R}(\theta) \begin{pmatrix} x \\ y \end{pmatrix} = \begin{pmatrix} \cos\theta & -\sin\theta \\ \sin\theta & \cos\theta \end{pmatrix} \begin{pmatrix} x \\ y \end{pmatrix} = \begin{pmatrix} x \cos\theta - y \sin\theta \\ x \sin\theta + y \cos\theta \end{pmatrix}
\end{equation}

By coupling these channels into a complex-valued scalar $Z = x + i y \in \mathbb{C}$, the action of a complex phase rotation by $\theta$ is represented by the Unitary group $U(1)$ via multiplication with the complex phasor $e^{i\theta}$:
\begin{align}
    Z' &= Z \cdot e^{i\theta} \\
       &= (x + i y)(\cos\theta + i \sin\theta) \\
       &= (x \cos\theta - y \sin\theta) + i(x \sin\theta + y \cos\theta)
\end{align}
Mapping the real and imaginary parts of $Z'$ back to the channel coordinates:
\begin{equation}
    \text{Re}(Z') = x \cos\theta - y \sin\theta, \quad \text{Im}(Z') = x \sin\theta + y \cos\theta
\end{equation}
This establishes an exact isomorphism between real-space 2D coordinate rotations $\mathbf{R}(\theta) \in SO(2)$ and complex-field phase multiplications $e^{i\theta} \in U(1)$. By utilizing this isomorphism, QSPA models and manipulates rotational coordinate dynamics directly in the complex domain.

\subsection{Preservation of Representational Similarity Metrics}
Let $\mathbf{z}_1, \mathbf{z}_2 \in \mathbb{C}^d$ denote the complex-valued activation vectors of two distinct data samples at layer $l$, where $d = C/2$. The complex-valued inner product between these samples is defined as:
\begin{equation}
    \langle \mathbf{z}_1, \mathbf{z}_2 \rangle_{\mathbb{C}} = \mathbf{z}_1^\dagger \mathbf{z}_2 = \sum_{j=1}^d Z_{1, j} \overline{Z_{2, j}}
\end{equation}
where $\dagger$ denotes the conjugate transpose and $\overline{Z}$ represents the complex conjugate.

Applying a channel-wise phase rotation vector $\boldsymbol{\theta} = (\theta_1, \dots, \theta_d)^T$ yields the transformed activations:
\begin{equation}
    Z'_{1, j} = Z_{1, j} e^{i\theta_j}, \quad Z'_{2, j} = Z_{2, j} e^{i\theta_j}
\end{equation}
The inner product of the transformed activations is given by:
\begin{align}
    \langle \mathbf{z}'_1, \mathbf{z}'_2 \rangle_{\mathbb{C}} &= \sum_{j=1}^d Z'_{1, j} \overline{Z'_{2, j}} \\
    &= \sum_{j=1}^d (Z_{1, j} e^{i\theta_j}) \left( \overline{Z_{2, j} e^{i\theta_j}} \right) \\
    &= \sum_{j=1}^d Z_{1, j} e^{i\theta_j} \overline{Z_{2, j}} e^{-i\theta_j} \\
    &= \sum_{j=1}^d Z_{1, j} \overline{Z_{2, j}} = \langle \mathbf{z}_1, \mathbf{z}_2 \rangle_{\mathbb{C}}
\end{align}
This proof demonstrates that the complex phase rotation is a unitary transformation that perfectly preserves the complex inner products between all samples. Since similarity indices such as cosine similarity, representational alignment, and Centered Kernel Alignment (CKA) similarity are functions of the inner product matrix, \textit{unitary phase rotations preserve the entire geometric and representational structure of the latent manifold}.

\subsection{Geometric Distortion under Independent Real Scaling}
In contrast, consider standard independent channel-wise scaling (e.g., TAAC), which scales channels $2j$ and $2j+1$ by independent real-valued scaling factors $s_{2j}$ and $s_{2j+1}$ estimated from a calibration set. The resulting complex activation components are:
\begin{equation}
    Z''_{1, j} = s_{2j} x_{1, j} + i s_{2j+1} y_{1, j}, \quad Z''_{2, j} = s_{2j} x_{2, j} + i s_{2j+1} y_{2, j}
\end{equation}
The complex inner product under this transformation is:
\begin{align}
    \langle \mathbf{z}''_1, \mathbf{z}''_2 \rangle_{\mathbb{C}} &= \sum_{j=1}^d Z''_{1, j} \overline{Z''_{2, j}} \\
    &= \sum_{j=1}^d \left( s_{2j} x_{1, j} + i s_{2j+1} y_{1, j} \right) \left( s_{2j} x_{2, j} - i s_{2j+1} y_{2, j} \right) \\
    &= \sum_{j=1}^d \left[ s_{2j}^2 x_{1, j} x_{2, j} + s_{2j+1}^2 y_{1, j} y_{2, j} + i s_{2j} s_{2j+1} \left( y_{1, j} x_{2, j} - x_{1, j} y_{2, j} \right) \right]
\end{align}
If $s_{2j} \neq s_{2j+1}$ (which is almost always the case when scaling parameters are estimated independently on finite, noisy calibration sets), the inner product is non-linearly distorted. Specifically, the relative angles, lengths, and orthogonalities between sample representations are warped. When compounded across multiple layers, this geometric warping destroys the semantic manifold structure of the representation space, contributing to representation collapse. 

By coupling adjacent channels into complex-valued wavefields, QSPA applies a single unified rotation $\theta_j$ and a joint amplitude scale $s^{\text{amp}}_j$ (which is highly stabilized via Tikhonov regularization $\lambda_{\text{reg}}$):
\begin{equation}
    s^{\text{amp}}_{l, j} = \sqrt{\frac{V^{(T)}_{l, j} + \lambda_{\text{reg}}}{V^{(M)}_{l, j} + \lambda_{\text{reg}}}}
\end{equation}
This unified treatment of channel pairs maintains coordinate harmony, prevents the unbalancing of adjacent features, and robustly resolves the Sparsity Trap under small calibration budgets.

\section{Experimental Hyperparameters and Reproducibility Details}
\label{appendix:reproducibility}

To guarantee 100\% scientific reproducibility, we document the precise hyperparameters, training regimes, and architectural configurations of our empirical campaign.

\subsection{Backbone Architecture (ResNet-18)}
We utilize a standard ResNet-18 backbone. The model processes $32 \times 32$ RGB images through the following sequential blocks:
\begin{enumerate}
    \item \textbf{Input Layer:} $3 \times 3$ Convolution (64 filters, stride 1, padding 1), Batch Normalization, and ReLU activation.
    \item \textbf{Layer 1:} Two residual blocks, each containing two $3\times3$ convolutions (64 filters) with batch normalization and skip connections.
    \item \textbf{Layer 2:} Two residual blocks (128 filters, downsampling via stride 2 on the first convolution).
    \item \textbf{Layer 3:} Two residual blocks (256 filters, downsampling via stride 2 on the first convolution).
    \item \textbf{Layer 4:} Two residual blocks (512 filters, downsampling via stride 2 on the first convolution).
    \item \textbf{Pooling and Classification:} Global Average Pooling (reducing the tensor from $512 \times 4 \times 4$ to $512$), followed by task-specific fully connected heads ($512 \times 10$).
\end{enumerate}
The model has a total of 20 BatchNorm2d layers, which are the targets of our full-backbone calibration sweeps.

\subsection{Expert Pre-Training Configurations}
Each expert model is trained independently on its respective task. Grayscale datasets (MNIST and Fashion-MNIST) are dynamically resized to $32 \times 32$ pixels and replicated across 3 channels. The training hyperparameters are specified in \cref{tab:hyperparameters}.

\begin{table}[H]
\caption{\textbf{Pre-training Hyperparameters} for MNIST, Fashion-MNIST, and CIFAR-10 experts.}
\label{tab:hyperparameters}
\begin{center}
\begin{tabular}{lc}
\toprule
\textbf{Hyperparameter} & \textbf{Value} \\
\midrule
Pre-trained Initialization & ImageNet-1k (via torchvision) \\
Optimizer & AdamW \\
Learning Rate & $5 \times 10^{-4}$ \\
Weight Decay & $10^{-4}$ \\
Learning Rate Schedule & Cosine Annealing \\
Batch Size & 64 \\
Training Epochs & 5 \\
Training Set Size (per task) & 3,000 samples (random subset) \\
Validation Set Size & 10,000 samples (full test split) \\
\bottomrule
\end{tabular}
\end{center}
\end{table}

This training regime produces expert models with stable test performance: \textbf{98.28\%} on MNIST, \textbf{87.64\%} on Fashion-MNIST, and \textbf{64.58\%} on CIFAR-10.

\subsection{Head Adaptation (REDA-style) Details}
For linear head-only adaptation sweeps ($N=256$ labeled samples per task, frozen ResNet backbone), we use:
\begin{itemize}
    \item \textbf{Optimizer:} Stochastic Gradient Descent (SGD) with momentum $0.9$.
    \item \textbf{Learning Rate:} $10^{-3}$ with constant schedule.
    \item \textbf{Batch Size:} 32.
    \item \textbf{Epochs:} 10.
    \item \textbf{Data Augmentation:} Standard random horizontal flips and crops for CIFAR-10; none for MNIST/Fashion-MNIST.
\end{itemize}

\section{Clifford-Valued Representation Alignment: The Multi-Dimensional Frontier}
\label{appendix:clifford}

In this appendix, we outline the theoretical generalization of Quantum Superposition Phase Alignment (QSPA) to higher-dimensional manifolds. While QSPA couples channels in pairs to exploit the isomorphism between $SO(2)$ and the complex phase group $U(1)$, real-world deep neural networks often exhibit higher-dimensional coordinate structures across channels. We introduce \textbf{Clifford-Valued Representation Alignment (CVRA)}, which utilizes Clifford (Geometric) Algebras to model multi-channel representations as hyper-dimensional wavefields, enabling geometry-preserving alignment across arbitrary dimensions.

\subsection{Clifford Algebras and Geometric Wavefields}
Let $d = 2^k$ represent the dimension of a grouped channel subspace. We construct the Clifford algebra $Cl_{0, d}(\mathbb{R})$ over the real vector space $\mathbb{R}^d$ equipped with a negative-definite quadratic form. The algebra is generated by $d$ orthogonal basis vectors $\{e_1, e_2, \dots, e_d\}$ satisfying the fundamental anti-commutation relations:
\begin{equation}
    e_i^2 = -1 \quad (\forall i), \quad e_i e_j = -e_j e_i \quad (\forall i \neq j)
\end{equation}
A general element in $Cl_{0, d}(\mathbb{R})$, called a multi-vector, is a linear combination of the scalar, vector, bivector, trivector, and higher-grade pseudoscalar elements:
\begin{equation}
    X = x_0 + \sum_{i=1}^d x_i e_i + \sum_{i < j} x_{ij} e_i e_j + \dots + x_{1 \dots d} e_1 e_2 \dots e_d
\end{equation}
To model high-dimensional activations, we group $d$ adjacent real-valued channels of layer $l$ into a Clifford multi-vector field. For each subspace $j \in \{0, \dots, C/d - 1\}$, the activation components are mapped to the vector part of a Clifford multi-vector:
\begin{equation}
    \mathbf{X}_{l, j} = \sum_{a=1}^d X_{l, d \cdot j + a} \, e_a
\end{equation}
This representation embeds the $d$-dimensional representation space directly into the geometric manifold of $Cl_{0, d}(\mathbb{R})$.

\subsection{The Spin Group and Rotors as Geometry-Preserving Operators}
In Clifford algebra, arbitrary $d$-dimensional rotations are represented elegantly by the Spin group $Spin(d)$, which is the double cover of the special orthogonal group $SO(d)$. An element $R \in Spin(d)$, called a \textit{rotor}, can be expressed as the geometric product of an even number of unit vectors and satisfies:
\begin{equation}
    R \widetilde{R} = 1
\end{equation}
where $\widetilde{R}$ denotes the reversion of $R$, which reverses the order of all basis vector products.

The rotation of a Clifford multi-vector activation $\mathbf{X}$ by a rotor $R$ is given by the sandwich product:
\begin{equation}
    \mathbf{X}' = R \mathbf{X} \widetilde{R}
\end{equation}
We now prove that this rotor action is a geometry-preserving unitary operation that preserves representational similarity. The geometric inner product between two multi-vectors $\mathbf{X}$ and $\mathbf{Y}$ is defined as the scalar part of their geometric product:
\begin{equation}
    \langle \mathbf{X}, \mathbf{Y} \rangle_{Cl} = \langle \mathbf{X} \widetilde{Y} \rangle_0 = \sum_{A} X_A Y_A
\end{equation}
Applying the rotor transformation to both activations:
\begin{align}
    \langle \mathbf{X}', \mathbf{Y}' \rangle_{Cl} &= \langle \left( R \mathbf{X} \widetilde{R} \right) \left( \widetilde{R \mathbf{Y} \widetilde{R}} \right) \rangle_0 \\
    &= \langle R \mathbf{X} \widetilde{R} R \mathbf{Y} \widetilde{R} \rangle_0
\end{align}
Using the rotor property $\widetilde{R} R = 1$:
\begin{align}
    \langle \mathbf{X}', \mathbf{Y}' \rangle_{Cl} &= \langle R \mathbf{X} \mathbf{Y} \widetilde{R} \rangle_0
\end{align}
By the cyclic property of the scalar part operator $\langle \cdot \rangle_0$:
\begin{align}
    \langle \mathbf{X}', \mathbf{Y}' \rangle_{Cl} &= \langle \widetilde{R} R \mathbf{X} \mathbf{Y} \rangle_0 \\
    &= \langle \mathbf{X} \mathbf{Y} \rangle_0 = \langle \mathbf{X}, \mathbf{Y} \rangle_{Cl}
\end{align}
This proof mathematically guarantees that \textit{Clifford rotor alignment is an exact, multidimensional isometric transformation}. It preserves the absolute angles, magnitudes, and similarity structures (including multi-dimensional CKA) of the latent activation manifold, establishing a completely non-warping foundation for multi-channel calibration.

\subsection{Multi-Task Consensus Clifford Rotors}
To align the merged model's representations on a joint calibration set $\mathcal{D}_{\text{joint}}$, we analyze the spatial covariance tensor in Clifford space. For each $d$-dimensional subspace, we define the Clifford covariance tensor $\mathbf{\Sigma} \in Cl_{0, d}(\mathbb{R}) \otimes Cl_{0, d}(\mathbb{R})$. 

The consensus orientation of the experts is represented by a target Clifford frame $\mathbf{F}^{(T)}$, computed as the Frechet mean of the experts' individual covariance frames under the Riemannian metric of the $Spin(d)$ manifold:
\begin{equation}
    \mathbf{F}^{(T)} = \arg\min_{\mathbf{F} \in Spin(d)} \sum_{k=1}^K \text{dist}^2_{Spin(d)}\left(\mathbf{F}, \mathbf{F}^{(k)}\right)
\end{equation}
Similarly, we extract the covariance frame of the merged model $\mathbf{F}^{(M)}$. The optimal alignment rotor $R_{l, j}$ is the unique rotor that maps $\mathbf{F}^{(M)}$ to the target consensus frame $\mathbf{F}^{(T)}$:
\begin{equation}
    R_{l, j} = \mathbf{F}^{(T)} \left(\mathbf{F}^{(M)}\right)^{-1}
\end{equation}
To regularize this alignment, the rotor can be interpolated toward the identity rotor using the Lie algebra generator (bivector) proportional to the multi-dimensional anisotropy ratio of the distribution, ensuring absolute robustness against sparse or isotropic channel noise.

\subsection{Blueprint for Natively Clifford-Valued Neural Networks}
While CVRA provides a post-hoc calibration mechanism, the ultimate realization of this visionary paradigm lies in constructing \textbf{Natively Clifford-Valued Neural Networks (CVNNs)}.

In a CVNN, the weights $\mathbf{W}$, biases $\mathbf{B}$, and activations $\mathbf{X}$ are all multi-vectors in $Cl_{0, d}(\mathbb{R})$. The forward pass of a Clifford-valued linear layer is defined using the Clifford geometric product:
\begin{equation}
    \mathbf{Y} = \mathbf{W} \mathbf{X} + \mathbf{B}
\end{equation}
Because the geometric product incorporates both the symmetric inner product and anti-symmetric outer product, a single Clifford layer inherently models multi-dimensional scaling, rotations, and projections. 

\textbf{Intrinsic Robustness to Representation Collapse:}
When merging two CVNN experts, the weights are averaged in Clifford space: $\mathbf{W}^{(M)} = \frac{1}{2}(\mathbf{W}^{(1)} + \mathbf{W}^{(2)})$. Because the geometric multiplication is equivariant under the action of the Spin group:
\begin{equation}
    \mathbf{W} \left( R \mathbf{X} \widetilde{R} \right) = R \left( \mathbf{W} \mathbf{X} \right) \widetilde{R}
\end{equation}
any coordinate drift between experts is represented as a rotor rotation in the Clifford manifold. Linear interpolation of Clifford weights corresponds to geodesic interpolation on the Spin group, which preserves the rotational and structural relationships between channels. Consequently, CVNNs are structurally immune to the representation collapse and coordinate scrambling that plague standard real-valued neural networks, paving the way for seamless, zero-shot multi-task model composition at scale.

\end{document}
